\scanpage{173}
\[
 \left\{\begin{aligned}
 C(x^{p-1}dx)&=dx\\
 C(dx/x)&=dx/x
 \end{aligned}\right.
\]
C'est le cas envisagé dans la note initiale de Cartier.

\textcircled{4} Fonctorialités, Globalisation.

Soit
\diagram{\begin{tikzcd}[ampersand replacement=\&,row sep=.65em]
 X\arrow[d]\\ S
\end{tikzcd}}
en car $p$, d'où
\diagram{\begin{tikzcd}[ampersand replacement=\&,column sep=2em]
 X\arrow[d,dash] \& X^{(p/S)}\arrow[l]\arrow[d,dash] \& X\arrow[l,"{f=f_{X/S}}"']\arrow[ll,bend right=25,"F_X"']\arrow[dl,dash]\\
 S \& S\arrow[l,"F_S"]\arrow[ul,phantom,"\text{\uline{cart.}}",very near start] \& {}
\end{tikzcd}}
, avec $f=f_{X/S}$ radiciel \deleted{\unsure{entier}} surjectif.

\deleted{Supposons que $f_*(\underline O_X)$ ait \unsure{une} localement une $p$-base sur $\operatorname{Im}(\underline O_{X^{(p/S)}}\longrightarrow f_*(\underline O_X))$, alors on obtient des isom. can.}

On définit alors un hom canonique de faisceaux d'Algèbres
\[
 \psi:\Omega^*_{X^{p/S}/S}\longrightarrow f_*(Z^*(\Omega^*_{X/S}))
\]
d'où un hom composé
\[
 \psi:\Omega^*_{X^{p/S}/S}\longrightarrow f_*(\mathscr H^*(\Omega^*_{X/S}))
 \simeq\mathscr H^\alpha(f_*(\deleted{\unsure{Z}}\!(\Omega^*_{X/S}))),
\]
Si $\underline\Lambda=\operatorname{Im}(\underline O_{X^{p/S}}\longrightarrow f_*\underline O_X))$ [qui est donc de la forme $\underline O_{X^{p/S}}/J$, $J$ un nilidéal, donc $J=0$ si $X^{p/S}$ est réduit, p.ex. $S$ réduit et $X$ \deleted{\unsure{lisse}}\inserted{loc} de présentation finie et \unsure{ouvert|u. ouvert} sur $S$], on aura donc
\[
 \widetilde\psi:\Omega^\alpha_{X^{p/S}/S}\otimes_{\underline O_{X^{p/S}}}\underline\Lambda
 \longrightarrow\deleted{f_*}(\mathscr H^\alpha(f_*(\Omega^*_{X/S})))
\]
Si $f_*(\underline O_X)$ admet loc une $p$-base sur $\underline O_{X^{p/S}}$ (ou ce qui revient au même, sur $\Lambda$), alors l'hom précédent est un isom, et on peut définir, \deleted{\unsure{la ca}} l'isom inverse
\[
 \widetilde C:\mathscr H^\alpha(f_*(\Omega^\cdot_{X/S}))
 \longrightarrow\Omega^*_{X^{p/S}/S}\otimes_{\underline O_{X^{p/S}}}\Lambda
\]
\margin{définir $\widetilde C$ chaque fois que $\psi$ est un isom (sans qu'il y ait néc. des $p$-bases\ldots)}
d'où un hom
\[
 \widetilde C:Z^*(f_*(\Omega^*_{X/S}))\longrightarrow
 \Omega^*_{\deleted{\unsure{X/S}}X^{p/S}/S}\otimes_{\underline O_{X^{p/S}}}\Lambda
\]
% archive sheet number: 172

